% GENERATED FILE. Edit canonical lessons, then run npm run generate:books.
% canonical-source-hash: fnv1a64:4daad99684e7d3c6
% canonical-lessons: TA-C420-calavai, TA-C420-cevai, TA-C420-cumaitukki, TA-C420-kuli-wage, TA-C420-tuppuravu

\chapter{Laundry and Porters}
\label{ch:ta-laundry-and-porters}
\hlchaptermodality{420}

\begin{chapteropening}
\textbf{By the end of this chapter:} \emph{I can talk about laundry, a service, porters and their wages, and cleaning.}

Complete the last lesson of chapter 420: I can talk about laundry, a service, porters and their wages, and cleaning.
\end{chapteropening}

\section[\texorpdfstring{\ta{சலவை} (calavai)}{சலவை (calavai)}]{\ta{சலவை} (calavai) — laundry; washing}

\label{lesson:TA-C420-calavai}



\begin{quote}
Before the new one: say the Tamil for to relax, then the Tamil for hospitality.
\end{quote}

\subsection*{You'll want to know: \ta{சலவை}}
\textbf{\ta{சலவை}} — \emph{calavai} — \textquotedblleft{}laundry; washing\textquotedblright{}.

\begin{quote}
\ta{என்} \ta{சட்டைகளைச்} \ta{சலவைக்குக்} \ta{கொடுத்தேன்}.
\end{quote}

\begin{quote}
\emph{eṉ caṭṭaikaḷaic calavaikkuk koṭuttēṉ.}
\end{quote}

\begin{quote}
I gave my shirts to the laundry.
\end{quote}

\ta{சலவைத்} \ta{தொழிலாளி} is a washerman.

\begin{cousinweb}
\textbf{\ta{சலவை}} (\emph{calavai}) is native Dravidian (DEDR 2368), \textquotedblleft{}the washing or bleaching of cloth\textquotedblright{}.
\end{cousinweb}

\begin{grammarlens}[title={What you've built}]
One more word for getting around.
\end{grammarlens}

\subsection*{Your turn}
\begin{itemize}
\raggedright
  \item \emph{Say it:} \emph{calavai}
  \item \emph{Say it:} \emph{calavai}, once more
  \item \emph{Say it:} \emph{eṉ caṭṭaikaḷaic calavaikkuk koṭuttēṉ.} — I gave my shirts to the laundry.
  \item \emph{Recall:} say the Tamil for to relax, then the Tamil for hospitality, then say \emph{calavai} again
\end{itemize}

\begin{tcolorbox}[breakable,colback=teal!4,colframe=teal!35!black,title={Before you move on}]
What does \ta{சலவை} mean? (\textquotedblleft{}Laundry; washing\textquotedblright{}.) Say it once more.
\end{tcolorbox}
\section[\texorpdfstring{\ta{சேவை} (cēvai)}{சேவை (cēvai)}]{\ta{சேவை} (cēvai) — a service}

\label{lesson:TA-C420-cevai}



\begin{quote}
Before the new one: say the Tamil for hospitality, then the Tamil for laundry.
\end{quote}

\subsection*{You'll want to know: \ta{சேவை}}
\textbf{\ta{சேவை}} — \emph{cēvai} — \textquotedblleft{}a service\textquotedblright{}.

\begin{quote}
\ta{இந்த} \ta{விடுதியின்} \ta{சேவை} \ta{மிக} \ta{நன்றாக} \ta{இருந்தது}.
\end{quote}

\begin{quote}
\emph{inta viṭutiyiṉ cēvai mika naṉṟāka iruntatu.}
\end{quote}

\begin{quote}
The service at this lodge was very good.
\end{quote}

\ta{பேருந்துச்} \ta{சேவை} is a bus service.

\begin{cousinweb}
\textbf{\ta{சேவை}} (\emph{cēvai}) is borrowed from Sanskrit \emph{sevā}, \textquotedblleft{}attendance, service\textquotedblright{} (CDIAL 13595).
\end{cousinweb}

\begin{grammarlens}[title={What you've built}]
One more word for getting around.
\end{grammarlens}

\subsection*{Your turn}
\begin{itemize}
\raggedright
  \item \emph{Say it:} \emph{cēvai}
  \item \emph{Say it:} \emph{cēvai}, once more
  \item \emph{Say it:} \emph{inta viṭutiyiṉ cēvai mika naṉṟāka iruntatu.} — The service at this lodge was very good.
  \item \emph{Recall:} say the Tamil for hospitality, then the Tamil for laundry, then say \emph{cēvai} again
\end{itemize}

\begin{tcolorbox}[breakable,colback=teal!4,colframe=teal!35!black,title={Before you move on}]
What does \ta{சேவை} mean? (\textquotedblleft{}A service\textquotedblright{}.) Say it once more.
\end{tcolorbox}
\section[\texorpdfstring{\ta{சுமைதூக்கி} (cumaitūkki)}{சுமைதூக்கி (cumaitūkki)}]{\ta{சுமைதூக்கி} (cumaitūkki) — a porter}

\label{lesson:TA-C420-cumaitukki}



\begin{quote}
Before the new one: say the Tamil for laundry, then the Tamil for a service.
\end{quote}

\subsection*{You'll want to know: \ta{சுமைதூக்கி}}
\textbf{\ta{சுமைதூக்கி}} — \emph{cumaitūkki} — \textquotedblleft{}a porter\textquotedblright{}.

\begin{quote}
\ta{சுமைதூக்கி} \ta{எங்கள்} \ta{பெட்டிகளை} \ta{ரயிலில்} \ta{ஏற்றினார்}.
\end{quote}

\begin{quote}
\emph{cumaitūkki eṅkaḷ peṭṭikaḷai rayilil ēṟṟiṉār.}
\end{quote}

\begin{quote}
The porter loaded our boxes onto the train.
\end{quote}

\ta{சுமை}, a load, + \ta{தூக்கி}, a lifter.

\begin{cousinweb}
\textbf{\ta{சுமைதூக்கி}} (\emph{cumaitūkki}) is built from \ta{சுமை}, \textquotedblleft{}a load\textquotedblright{} (DEDR 2677), and \ta{தூக்கி}, \textquotedblleft{}one who lifts\textquotedblright{}, from \ta{தூக்கு} (DEDR 3376).
\end{cousinweb}

\begin{grammarlens}[title={What you've built}]
One more word for getting around.
\end{grammarlens}

\subsection*{Your turn}
\begin{itemize}
\raggedright
  \item \emph{Say it:} \emph{cumaitūkki}
  \item \emph{Say it:} \emph{cumaitūkki}, once more
  \item \emph{Say it:} \emph{cumaitūkki eṅkaḷ peṭṭikaḷai rayilil ēṟṟiṉār.} — The porter loaded our boxes onto the train.
  \item \emph{Recall:} say the Tamil for laundry, then the Tamil for a service, then say \emph{cumaitūkki} again
\end{itemize}

\begin{tcolorbox}[breakable,colback=teal!4,colframe=teal!35!black,title={Before you move on}]
What does \ta{சுமைதூக்கி} mean? (\textquotedblleft{}A porter\textquotedblright{}.) Say it once more.
\end{tcolorbox}
\section[\texorpdfstring{\ta{கூலி} (kūli)}{கூலி (kūli)}]{\ta{கூலி} (kūli) — a wage; a fee for labour}

\label{lesson:TA-C420-kuli-wage}



\begin{quote}
Before the new one: say the Tamil for a service, then the Tamil for a porter.
\end{quote}

\subsection*{You'll want to know: \ta{கூலி}}
\textbf{\ta{கூலி}} — \emph{kūli} — \textquotedblleft{}a wage; a fee for labour\textquotedblright{}.

\begin{quote}
\ta{சுமைதூக்கிக்கு} \ta{நூறு} \ta{ரூபாய்} \ta{கூலி} \ta{கொடுத்தோம்}.
\end{quote}

\begin{quote}
\emph{cumaitūkkikku nūṟu rūpāy kūli koṭuttōm.}
\end{quote}

\begin{quote}
We paid the porter a hundred rupees for his work.
\end{quote}

\ta{கூலித்} \ta{தொழிலாளி} is a daily-wage labourer.

\begin{cousinweb}
\textbf{\ta{கூலி}} (\emph{kūli}) is native Dravidian (DEDR 1905), \textquotedblleft{}wages, pay, hire\textquotedblright{}.
\end{cousinweb}

\begin{grammarlens}[title={What you've built}]
One more word for getting around.
\end{grammarlens}

\subsection*{Your turn}
\begin{itemize}
\raggedright
  \item \emph{Say it:} \emph{kūli}
  \item \emph{Say it:} \emph{kūli}, once more
  \item \emph{Say it:} \emph{cumaitūkkikku nūṟu rūpāy kūli koṭuttōm.} — We paid the porter a hundred rupees for his work.
  \item \emph{Recall:} say the Tamil for a service, then the Tamil for a porter, then say \emph{kūli} again
\end{itemize}

\begin{tcolorbox}[breakable,colback=teal!4,colframe=teal!35!black,title={Before you move on}]
What does \ta{கூலி} mean? (\textquotedblleft{}A wage; a fee for labour\textquotedblright{}.) Say it once more.
\end{tcolorbox}
\section[\texorpdfstring{\ta{துப்புரவு} (tuppuravu)}{துப்புரவு (tuppuravu)}]{\ta{துப்புரவு} (tuppuravu) — cleaning; cleanliness}

\label{lesson:TA-C420-tuppuravu}



\begin{quote}
Before the new one: say the Tamil for a porter, then the Tamil for a wage.
\end{quote}

\subsection*{You'll want to know: \ta{துப்புரவு}}
\textbf{\ta{துப்புரவு}} — \emph{tuppuravu} — \textquotedblleft{}cleaning; cleanliness\textquotedblright{}.

\begin{quote}
\ta{தினமும்} \ta{காலையில்} \ta{அறைகளில்} \ta{துப்புரவு} \ta{வேலை} \ta{நடக்கிறது}.
\end{quote}

\begin{quote}
\emph{tiṉamum kālaiyil aṟaikaḷil tuppuravu vēlai naṭakkiṟatu.}
\end{quote}

\begin{quote}
Every morning the rooms are cleaned.
\end{quote}

\ta{துப்புரவுப்} \ta{பணி} is cleaning work.

\begin{cousinweb}
\textbf{\ta{துப்புரவு}} (\emph{tuppuravu}), \textquotedblleft{}cleaning, cleanliness\textquotedblright{}, is a Tamil word whose origin is not settled.
\end{cousinweb}

\begin{grammarlens}[title={What you've built}]
One more word for getting around.
\end{grammarlens}

\subsection*{Your turn}
\begin{itemize}
\raggedright
  \item \emph{Say it:} \emph{tuppuravu}
  \item \emph{Say it:} \emph{tuppuravu}, once more
  \item \emph{Say it:} \emph{tiṉamum kālaiyil aṟaikaḷil tuppuravu vēlai naṭakkiṟatu.} — Every morning the rooms are cleaned.
  \item \emph{Recall:} say the Tamil for a porter, then the Tamil for a wage, then say \emph{tuppuravu} again
\end{itemize}

\begin{tcolorbox}[breakable,colback=teal!4,colframe=teal!35!black,title={Before you move on}]
What does \ta{துப்புரவு} mean? (\textquotedblleft{}Cleaning; cleanliness\textquotedblright{}.) Say it once more.
\end{tcolorbox}
